\documentclass[11pt]{article}
\usepackage{definitions}
\begin{document}

\textbf{\uppercase{Statement of Work --- Yale University}}\\[.5em]

For this project, the Yale team (PI: Dr.\ Martin Pfaller) will develop and validate a patient-specific constrained mixture model to predict ventricular growth following the ventricular recruitment procedure (Aim 3). Building on our homogenized constrained mixture framework for cardiac \gls*{gr}, we will extend the reduced-geometry ventricular model into a mixture model in which the myocardial wall comprises contractile and passive constituents whose deposition and degradation are driven by deviations of wall stress from a homeostatic set-point. Each model will be initialized from the pre-recruitment Digital Twin provided by the primary site and driven by the post-recruitment mechanical loading; we will assist in defining the growth stimulus and in calibrating constituent turnover and gain parameters, drawing patient-specific geometry, loading, and baseline wall stress from the invasive cardiac MRI data and adopting established priors for parameters that cannot be measured in vivo. These tools will be implemented in our open-source modeling framework and verified against published cardiac \gls*{gr} benchmark cases. After building the framework, we will apply it to the recruitment cohort, integrating each model forward over the 8--15 month post-recruitment interval to predict remodeling of chamber volume, mass, and wall thickness. Finally, we will validate these predictions against the paired post-recruitment exams and assess whether the model identifies which ventricles reach a size adequate for biventricular repair.

\end{document}
